% !TEX spellcheck = en_US



% ~~~~~~~~~~~~~~~~~~~~~~~~~~~~~~~~~~~~~~~ V
% (NC) 2026 Haluk Bingol
% github.com/halukbingol/LaTeX-Templates
%
% Licensees may copy, distribute, display, and perform the work and make derivative works 
% and remixes based on it only for non-commercial purposes. 
% ~~~~~~~~~~~~~~~~~~~~~~~~~~~~~~~~~~~~~~~ A




% v20261123T2200

\documentclass[aspectratio=169,10pt,compress]{beamer}




% ~~~~~~~~~~~~~~~~~~~~~~~~~~~~~~~~~~~~~~~~~~~~~~~~~~~~~~~~~~~~~~~~~ V hblib
	\usepackage{../../hblib/hblibPresentation} % lib presentation: theme, i/N, \hSection, algorithm2e, ...
	\usepackage{../../hblib/hblibCodeoutput} % lib code + output: \lstcode, \lstterminal, ...
	\usepackage{../../hblib/hblibDatastructures} % lib data structures: \hString, \hStack, \hQueue, \hArray, \hGrid
% ~~~~~~~~~~~~~~~~~~~~~~~~~~~~~~~~~~~~~~~~~~~~~~~~~~~~~~~~~~~~~~~~~ A hblib




% +++++++++++++++++++++++++++++++++++++++ V specific
\DontPrintSemicolon % in algorithms
% +++++++++++++++++++++++++++++++++++++++ A specific



\title[Exceptions, Tests: Sophomore]{
	Exception Handling and Unit Testing\\
	Sophomore
}
\subtitle{
	\cisOOO\ -- Level 2: design, libraries, tools, good habits
}
\author[Erkovan]{Sait Erkovan}

\institute[FCIS]{
	Faculty of Computer and Information Sciences\\
	Yeditepe University
}
\date{{\tiny \hbTimeStamp}}

\begin{document}




% ~~~~~~~~~~~~~~~~~~~~~~~~~~~~~~~~~~~~~~ V frm
\begin{frame}[t,fragile]
	\titlepage
\end{frame}
% ~~~~~~~~~~~~~~~~~~~~~~~~~~~~~~~~~~~~~~ A frm




% ~~~~~~~~~~~~~~~~~~~~~~~~~~~~~~~~~~~~~~ V frm
\begin{frame}[t,fragile] \frametitle{Outline}
	\begin{columns}[T]
		\begin{column}{0.48\textwidth}
			\tableofcontents[sections={1-3}, hideallsubsections]
		\end{column}
		\begin{column}{0.48\textwidth}
			\tableofcontents[sections={4-7}, hideallsubsections]
		\end{column}
	\end{columns}

	\vspace{1em}
	\small It builds on \hCode{-1-freshman}: the messages, \hCode{try}/\hCode{catch}/\hCode{finally} and first tests are assumed.
	Classes with fields and constructors (CIS112) are used from here on.
\end{frame}
% ~~~~~~~~~~~~~~~~~~~~~~~~~~~~~~~~~~~~~~ A frm




% ~~~~~~~~~~~~~~~~~~~~~~~~~~~~~~~~~~~~~~ V
\hSection[Contract]{A class with a contract}
% ~~~~~~~~~~~~~~~~~~~~~~~~~~~~~~~~~~~~~~ A




% ~~~~~~~~~~~~~~~~~~~~~~~~~~~~~~~~~~~~~~ V frm
\begin{frame}[t,fragile] \frametitle{Unchecked: the caller's mistakes}
	\lstcode[basicstyle=\ttfamily\fontsize{6.4}{7.6}\selectfont, firstline=9, lastline=25, firstnumber=9]{codeoutput/Bank/Account.java}{Java}

	\footnotesize Fields \hCode{owner}, \hCode{balance}, \hCode{closed} (lines 5--7). A negative amount, a \hCode{null} owner, a closed
	account: bugs of the calling code. Unchecked, documented with \hCode{@throws}; \hCode{checkOpen} is at line 51.
\end{frame}
% ~~~~~~~~~~~~~~~~~~~~~~~~~~~~~~~~~~~~~~ A frm




% ~~~~~~~~~~~~~~~~~~~~~~~~~~~~~~~~~~~~~~ V frm
\begin{frame}[t,fragile] \frametitle{Checked: a normal problem the caller must handle}
	\lstcode[basicstyle=\ttfamily\fontsize{6.4}{7.6}\selectfont, firstline=27, lastline=41, firstnumber=27]{codeoutput/Bank/Account.java}{Java}

	\begin{itemize}\footnotesize
		\item Not enough money is not a bug: a correct program meets it every day. The caller must plan for it, so it is checked
		\item \hCode{throws InsufficientFundsException} is part of the method's signature, like its parameters
		\item The check comes \emph{before} the change (line 40): a failed call leaves the balance as it was
	\end{itemize}
\end{frame}
% ~~~~~~~~~~~~~~~~~~~~~~~~~~~~~~~~~~~~~~ A frm




% ~~~~~~~~~~~~~~~~~~~~~~~~~~~~~~~~~~~~~~ V frm
\begin{frame}[t,fragile] \frametitle{Standard exceptions to reuse}


% +++++++++++++++++++++++++++++++++++++++ V tbl
%: -Tbl. standard exceptions
\begin{table}
		\footnotesize
		\begin{tabular}{@{}L{0.34\textwidth}L{0.30\textwidth}L{0.30\textwidth}@{}}
			\toprule
			\textbf{Exception} & \textbf{When} & \textbf{Shortcut} \\
			\midrule
			\hCode{IllegalArgumentException}       & a parameter has a wrong value       & \\
			\hCode{NullPointerException}           & a parameter is \hCode{null}          & \hCode{Objects.requireNonNull(x, "x")} \\
			\hCode{IndexOutOfBoundsException}      & an index is out of range             & \hCode{Objects.checkIndex(i, n)} \\
			\hCode{IllegalStateException}          & the object is in the wrong state     & \\
			\hCode{NoSuchElementException}         & nothing left to return               & \\
			\hCode{UnsupportedOperationException}  & this object cannot do that at all    & \\
			\hCode{ArithmeticException}            & e.g. an \hCode{int} division by zero & \hCode{Math.addExact} (overflow) \\
			\bottomrule
		\end{tabular}
\end{table}%
% +++++++++++++++++++++++++++++++++++++++ A tbl


	\vspace{0.3em}
	\small Reuse these before writing your own: every Java programmer knows what they mean. All are unchecked.
\end{frame}
% ~~~~~~~~~~~~~~~~~~~~~~~~~~~~~~~~~~~~~~ A frm




% ~~~~~~~~~~~~~~~~~~~~~~~~~~~~~~~~~~~~~~ V frm
\begin{frame}[t,fragile] \frametitle{Using the account}
	\lstcode[basicstyle=\ttfamily\fontsize{6}{7.2}\selectfont, firstline=2, lastline=24, firstnumber=2]{codeoutput/Bank/AccountDemo.java}{Java}
	\only<2->{\begin{textblock*}{0.66\textwidth}(5.4cm,4.6cm)
		\lstoutput[basicstyle=\ttfamily\fontsize{6.4}{7.6}\selectfont, firstline=3, lastline=10]{codeoutput/Bank.txt}
	\end{textblock*}}
\end{frame}
% ~~~~~~~~~~~~~~~~~~~~~~~~~~~~~~~~~~~~~~ A frm




% ~~~~~~~~~~~~~~~~~~~~~~~~~~~~~~~~~~~~~~ V
\hSection[Own]{Your own exception types}
% ~~~~~~~~~~~~~~~~~~~~~~~~~~~~~~~~~~~~~~ A




% ~~~~~~~~~~~~~~~~~~~~~~~~~~~~~~~~~~~~~~ V frm
\begin{frame}[t,fragile] \frametitle{A checked exception of your own}
	\lstcode[basicstyle=\ttfamily\scriptsize]{codeoutput/Bank/InsufficientFundsException.java}{Java}

	\begin{itemize}\footnotesize
		\item \hCode{extends Exception}: checked; \hCode{extends RuntimeException}: unchecked
		\item \hCode{super(message)} keeps the message; a field adds details the caller can use (\hCode{getMissing})
		\item The name ends in \hCode{Exception} and says what went wrong, not where
	\end{itemize}
\end{frame}
% ~~~~~~~~~~~~~~~~~~~~~~~~~~~~~~~~~~~~~~ A frm




% ~~~~~~~~~~~~~~~~~~~~~~~~~~~~~~~~~~~~~~ V frm
\begin{frame}[t,fragile] \frametitle{Checked or unchecked? Own type or standard?}


% +++++++++++++++++++++++++++++++++++++++ V tbl
%: -Tbl. choosing the kind
\begin{table}
		\footnotesize
		\begin{tabular}{@{}L{0.40\textwidth}L{0.24\textwidth}L{0.28\textwidth}@{}}
			\toprule
			\textbf{The problem} & \textbf{Kind} & \textbf{Example} \\
			\midrule
			the caller broke the contract              & unchecked, standard & \hCode{IllegalArgumentException} \\
			a correct program meets it; callers can recover & checked          & \hCode{InsufficientFundsException} \\
			the outside world failed (file, network)    & checked, standard   & \hCode{IOException} \\
			callers need details or a special reaction  & your own type       & with fields such as \hCode{getMissing} \\
			nobody can recover                          & unchecked           & let it go up to a top-level handler \\
			\bottomrule
		\end{tabular}
\end{table}%
% +++++++++++++++++++++++++++++++++++++++ A tbl


	\vspace{0.3em}
	\small When in doubt, unchecked: a checked exception forces every caller, and every caller of those, to change.
\end{frame}
% ~~~~~~~~~~~~~~~~~~~~~~~~~~~~~~~~~~~~~~ A frm




% ~~~~~~~~~~~~~~~~~~~~~~~~~~~~~~~~~~~~~~ V
\hSection[Chain]{Translating and chaining}
% ~~~~~~~~~~~~~~~~~~~~~~~~~~~~~~~~~~~~~~ A




% ~~~~~~~~~~~~~~~~~~~~~~~~~~~~~~~~~~~~~~ V frm
\begin{frame}[t,fragile] \frametitle{Translate the exception, keep the cause}
	\lstcode[basicstyle=\ttfamily\fontsize{6}{7.2}\selectfont, firstline=3, lastline=25, firstnumber=3]{codeoutput/Records.java}{Java}
\end{frame}
% ~~~~~~~~~~~~~~~~~~~~~~~~~~~~~~~~~~~~~~ A frm




% ~~~~~~~~~~~~~~~~~~~~~~~~~~~~~~~~~~~~~~ V frm
\begin{frame}[t,fragile] \frametitle{Caused by}
	\lstoutput[basicstyle=\ttfamily\fontsize{6.4}{7.6}\selectfont]{codeoutput/Records.txt}

	\begin{itemize}\small
		\item First part: the new exception, in the words of this program: which record, which line
		\item \hCode{Caused by:} the original exception, with its own trace
		\item \hCode{... 1 more}: the rest of the trace is the same as above, so it is not repeated
		\item Without \hCode{e} in line 23, the cause would be lost: never translate without passing the cause
	\end{itemize}
\end{frame}
% ~~~~~~~~~~~~~~~~~~~~~~~~~~~~~~~~~~~~~~ A frm




% ~~~~~~~~~~~~~~~~~~~~~~~~~~~~~~~~~~~~~~ V
\hSection[Tools]{More language tools}
% ~~~~~~~~~~~~~~~~~~~~~~~~~~~~~~~~~~~~~~ A




% ~~~~~~~~~~~~~~~~~~~~~~~~~~~~~~~~~~~~~~ V frm
\begin{frame}[t,fragile] \frametitle{Multi-catch: one handler, several types}
	\lstcodedownstep[basicstyle=\ttfamily\fontsize{6.4}{7.6}\selectfont, firstline=2, lastline=15, firstnumber=2][basicstyle=\ttfamily\scriptsize]{codeoutput/MultiCatch.java}{Java}{codeoutput/MultiCatch.txt}

	\footnotesize \hCode{catch (A | B e)}: the same handling for both. The types may not be in one family (no parent with its child).
\end{frame}
% ~~~~~~~~~~~~~~~~~~~~~~~~~~~~~~~~~~~~~~ A frm




% ~~~~~~~~~~~~~~~~~~~~~~~~~~~~~~~~~~~~~~ V frm
\begin{frame}[t,fragile] \frametitle{Two resources in one try}
	\begin{columns}[T]
		\begin{column}{0.62\textwidth}
			\lstcode[basicstyle=\ttfamily\fontsize{6.4}{7.6}\selectfont, firstline=6, lastline=17, firstnumber=6]{codeoutput/CopyUpper/CopyUpper.java}{Java}
		\end{column}
		\begin{column}{0.36\textwidth}
			\lstcode[basicstyle=\ttfamily\scriptsize, numbers=none]{codeoutput/CopyUpper/names.txt}{}
			\uncover<2->{\lstterminal[basicstyle=\ttfamily\scriptsize]{codeoutput/CopyUpper.txt}}
		\end{column}
	\end{columns}

	\footnotesize Separated by \hCode{;} in the parentheses. Closed in reverse order: \hCode{out} first, so everything written is flushed,
	then \hCode{in}. No \hCode{finally}, no \hCode{close()} by hand.
\end{frame}
% ~~~~~~~~~~~~~~~~~~~~~~~~~~~~~~~~~~~~~~ A frm




% ~~~~~~~~~~~~~~~~~~~~~~~~~~~~~~~~~~~~~~ V frm
\begin{frame}[t,fragile] \frametitle{Log instead of print}
	\lstcode[basicstyle=\ttfamily\fontsize{6.4}{7.6}\selectfont, firstline=11, lastline=17, firstnumber=11]{codeoutput/LogDemo.java}{Java}
	\uncover<2->{\lstoutput[basicstyle=\ttfamily\fontsize{6}{7.2}\selectfont, lastline=8]{codeoutput/LogDemo.txt}}

	\footnotesize \uncover<2->{Levels (\hCode{SEVERE}, \hCode{WARNING}, \hCode{INFO}) can be switched on and off without changing the code;
	with \hCode{e} (line 16) the whole trace is written. Lines 9--10 (not shown) only set the format for this slide.}
\end{frame}
% ~~~~~~~~~~~~~~~~~~~~~~~~~~~~~~~~~~~~~~ A frm




% ~~~~~~~~~~~~~~~~~~~~~~~~~~~~~~~~~~~~~~ V
\hSection[JUnit]{JUnit in depth}
% ~~~~~~~~~~~~~~~~~~~~~~~~~~~~~~~~~~~~~~ A




% ~~~~~~~~~~~~~~~~~~~~~~~~~~~~~~~~~~~~~~ V frm
\begin{frame}[t,fragile] \frametitle{A fixture: @BeforeEach}
	\lstcode[basicstyle=\ttfamily\fontsize{6.4}{7.6}\selectfont, firstline=9, lastline=23, firstnumber=9]{codeoutput/Bank/AccountTest.java}{Java}

	\footnotesize Every test starts from the same \hDefined{fixture}: an account of Ada with 100. JUnit makes a new \hCode{AccountTest}
	object for every test, then calls \hCode{setUp}: no test sees what another one changed.
\end{frame}
% ~~~~~~~~~~~~~~~~~~~~~~~~~~~~~~~~~~~~~~ A frm




% ~~~~~~~~~~~~~~~~~~~~~~~~~~~~~~~~~~~~~~ V frm
\begin{frame}[t,fragile] \frametitle{Exceptions: the type, the message, the details}
	\lstcode[basicstyle=\ttfamily\fontsize{6.4}{7.6}\selectfont, firstline=32, lastline=49, firstnumber=32]{codeoutput/Bank/AccountTest.java}{Java}
	\only<2->{\begin{textblock*}{0.56\textwidth}(7.4cm,5.6cm)
		\lstoutput[basicstyle=\ttfamily\fontsize{6.4}{7.6}\selectfont, firstline=11]{codeoutput/Bank.txt}
	\end{textblock*}}
\end{frame}
% ~~~~~~~~~~~~~~~~~~~~~~~~~~~~~~~~~~~~~~ A frm




% ~~~~~~~~~~~~~~~~~~~~~~~~~~~~~~~~~~~~~~ V frm
\begin{frame}[t,fragile] \frametitle{What the last frame used}
	\begin{itemize}\small
		\item \hCode{assertThrows} returns the exception: check its message and its fields (lines 35, 46)
		\item \hCode{assertAll} (line 37): runs all its checks and reports every failure, not only the first
		\item \hCode{@DisplayName}: a sentence in the report instead of the method name
		\item A test may declare \hCode{throws} (line 27): a checked exception that is not expected simply makes the test an error
		\item \hCode{closed} (lines 51--56) checks the \hCode{IllegalStateException} of a closed account
		\item The state after a failed call is tested too (line 40): the balance did not change
	\end{itemize}
\end{frame}
% ~~~~~~~~~~~~~~~~~~~~~~~~~~~~~~~~~~~~~~ A frm




% ~~~~~~~~~~~~~~~~~~~~~~~~~~~~~~~~~~~~~~ V frm
\begin{frame}[t,fragile] \frametitle{The lifecycle of a test class}
	\begin{columns}[T]
		\begin{column}{0.56\textwidth}
			\lstcode[basicstyle=\ttfamily\fontsize{6}{7.2}\selectfont, firstline=7, lastline=22, firstnumber=7]{codeoutput/JUnitMore/LifecycleTest.java}{Java}
		\end{column}
		\begin{column}{0.42\textwidth}
			\uncover<2->{\lstoutput[basicstyle=\ttfamily\fontsize{6.4}{7.6}\selectfont, firstline=3]{codeoutput/Lifecycle.txt}}
		\end{column}
	\end{columns}

	\footnotesize \hCode{@AfterEach} and \hCode{@AfterAll} (lines 28--36) print the closing lines. \hCode{@BeforeAll} and
	\hCode{@AfterAll} are \hCode{static}: once per class, for expensive set-up such as a database.
\end{frame}
% ~~~~~~~~~~~~~~~~~~~~~~~~~~~~~~~~~~~~~~ A frm




% ~~~~~~~~~~~~~~~~~~~~~~~~~~~~~~~~~~~~~~ V frm
\begin{frame}[t,fragile] \frametitle{assertAll reports every failure}
	\lstcode[basicstyle=\ttfamily\scriptsize, firstline=6, lastline=19, firstnumber=6]{codeoutput/JUnitMore/AllTest.java}{Java}

	\small \hQuestion{Both tests expect A for 89 and for 80 (wrong on purpose). How many failures does each one report?}
\end{frame}
% ~~~~~~~~~~~~~~~~~~~~~~~~~~~~~~~~~~~~~~ A frm




% ~~~~~~~~~~~~~~~~~~~~~~~~~~~~~~~~~~~~~~ V frm
\begin{frame}[t,fragile] \frametitle{assertAll: the report}
	\lstoutput[basicstyle=\ttfamily\scriptsize, firstline=3]{codeoutput/All.txt}

	\begin{itemize}\small
		\item \hCode{separate} stops at the first false assert: the problem at 80 stays hidden until 89 is fixed
		\item \hCode{together} runs every check and lists all failures under its heading
		\item Use it for several properties of one result; for many inputs, a parameterized test is better
	\end{itemize}
\end{frame}
% ~~~~~~~~~~~~~~~~~~~~~~~~~~~~~~~~~~~~~~ A frm




% ~~~~~~~~~~~~~~~~~~~~~~~~~~~~~~~~~~~~~~ V frm
\begin{frame}[t,fragile] \frametitle{Parameterized tests: one method, a table of cases}
	\begin{columns}[T]
		\begin{column}{0.64\textwidth}
			\lstcode[basicstyle=\ttfamily\fontsize{6}{7.2}\selectfont, firstline=8, lastline=21, firstnumber=8]{codeoutput/JUnitMore/GradesParamTest.java}{Java}
		\end{column}
		\begin{column}{0.34\textwidth}
			\uncover<2->{\lstoutput[basicstyle=\ttfamily\fontsize{6}{7.2}\selectfont, firstline=3]{codeoutput/Param.txt}}
		\end{column}
	\end{columns}

	\footnotesize Each row of \hCode{@CsvSource} is one test; JUnit converts \hCode{"60"} and \hCode{"D"} to \hCode{int} and \hCode{char}.
	Needs \hCode{junit-jupiter-params} (in Eclipse's JUnit 5 library).
\end{frame}
% ~~~~~~~~~~~~~~~~~~~~~~~~~~~~~~~~~~~~~~ A frm




% ~~~~~~~~~~~~~~~~~~~~~~~~~~~~~~~~~~~~~~ V frm
\begin{frame}[t,fragile] \frametitle{@Timeout and @Disabled}
	\lstcode[basicstyle=\ttfamily\fontsize{6.4}{7.6}\selectfont, firstline=12, lastline=20, firstnumber=12]{codeoutput/JUnitMore/TimeTest.java}{Java}
	\lstcode[basicstyle=\ttfamily\fontsize{6.4}{7.6}\selectfont, firstline=27, lastline=36, firstnumber=27]{codeoutput/JUnitMore/TimeTest.java}{Java}
\end{frame}
% ~~~~~~~~~~~~~~~~~~~~~~~~~~~~~~~~~~~~~~ A frm




% ~~~~~~~~~~~~~~~~~~~~~~~~~~~~~~~~~~~~~~ V frm
\begin{frame}[t,fragile] \frametitle{What the timeout did}
	\lstoutput[basicstyle=\ttfamily\fontsize{6.4}{7.6}\selectfont, firstline=3]{codeoutput/Time.txt}

	\begin{itemize}\small
		\item \hCode{halvings(0)}: $0 / 2$ is 0 forever, so the loop never ends; without a limit, the whole test run would hang
		\item \hCode{@Timeout} stops waiting after 500 ms; the test is an \hDefined{error} (a \hCode{TimeoutException}), the others go on
		\item \hCode{SEPARATE\_THREAD} runs the test in its own thread, so a loop that never checks for interruption can be left behind
		\item \hCode{@Disabled("reason")}: skipped, but listed with the reason; better than deleting or commenting out a test
		\item The course grader also limits the time of every test: an endless loop costs that test, not the others
	\end{itemize}
\end{frame}
% ~~~~~~~~~~~~~~~~~~~~~~~~~~~~~~~~~~~~~~ A frm




% ~~~~~~~~~~~~~~~~~~~~~~~~~~~~~~~~~~~~~~ V
\hSection[Habits]{Good habits and tools}
% ~~~~~~~~~~~~~~~~~~~~~~~~~~~~~~~~~~~~~~ A




% ~~~~~~~~~~~~~~~~~~~~~~~~~~~~~~~~~~~~~~ V frm
\begin{frame}[t,fragile] \frametitle{Shaping a test}


% +++++++++++++++++++++++++++++++++++++++ V tbl
%: -Tbl. test habits
\begin{table}
		\footnotesize
		\begin{tabular}{@{}L{0.26\textwidth}L{0.66\textwidth}@{}}
			\toprule
			\textbf{Habit} & \textbf{Because} \\
			\midrule
			arrange, act, assert         & every reader finds the input, the call and the check at once \\
			one behavior per test        & a red test then names exactly one broken behavior \\
			a name that is a sentence    & the report reads as a list of what works and what does not \\
			no \hCode{if} or loops in a test & a test with logic needs its own tests; use a parameterized test \\
			expected values by hand      & a value copied from the program's output tests nothing \\
			test the exceptions too      & the \hCode{@throws} lines of the contract are behavior as well \\
			independent and repeatable   & any order, any machine, any day: no shared state, no clock, no random \\
			\bottomrule
		\end{tabular}
\end{table}%
% +++++++++++++++++++++++++++++++++++++++ A tbl
\end{frame}
% ~~~~~~~~~~~~~~~~~~~~~~~~~~~~~~~~~~~~~~ A frm




% ~~~~~~~~~~~~~~~~~~~~~~~~~~~~~~~~~~~~~~ V frm
\begin{frame}[t,fragile] \frametitle{Where the tests live}
	\begin{columns}[T]
		\begin{column}{0.48\textwidth}
			\small \textbf{Eclipse project}


% +++++++++++++++++++++++++++++++++++++++ V lst
%: -Lst. eclipse layout
\begin{lstlisting}[style=codesnippet, language={}, basicstyle=\ttfamily\scriptsize]
Bank/
  src/                  source folder
    bank/Account.java
  test/                 a second source folder
    bank/AccountTest.java
  JUnit 5               library on the build path
\end{lstlisting}
% +++++++++++++++++++++++++++++++++++++++ A lst


		\end{column}
		\begin{column}{0.48\textwidth}
			\small \textbf{Maven or Gradle project}


% +++++++++++++++++++++++++++++++++++++++ V lst
%: -Lst. maven layout
\begin{lstlisting}[style=codesnippet, language={}, basicstyle=\ttfamily\scriptsize]
bank/
  pom.xml  (or build.gradle)
  src/main/java/bank/Account.java
  src/test/java/bank/AccountTest.java
\end{lstlisting}
% +++++++++++++++++++++++++++++++++++++++ A lst


		\end{column}
	\end{columns}

	\vspace{0.5em}
	\begin{itemize}\small
		\item Tests apart from the code: the program can be shipped without them
		\item The same package (\hCode{bank}) in both: a test may use what is not \hCode{private}
		\item In Eclipse: \menu{New > Source Folder}, name \hCode{test}; the JUnit wizard asks which source folder to use
	\end{itemize}
\end{frame}
% ~~~~~~~~~~~~~~~~~~~~~~~~~~~~~~~~~~~~~~ A frm




% ~~~~~~~~~~~~~~~~~~~~~~~~~~~~~~~~~~~~~~ V frm
\begin{frame}[t,fragile] \frametitle{Running tests outside Eclipse}
	\small JUnit's Console Launcher, one jar from Maven Central (\hCode{junit-platform-console-standalone}): \hRemark{Not run here}


% +++++++++++++++++++++++++++++++++++++++ V lst
%: -Lst. console launcher
\begin{lstlisting}[style=codesnippet, language=bash, basicstyle=\ttfamily\scriptsize]
# Windows (Command Prompt):  ; between class path entries
javac -d out -cp junit-platform-console-standalone-6.1.3.jar src\*.java test\*.java
java -jar junit-platform-console-standalone-6.1.3.jar execute -cp out --scan-class-path
# macOS (Terminal):  : between class path entries, / in paths
javac -d out -cp junit-platform-console-standalone-6.1.3.jar src/*.java test/*.java
java -jar junit-platform-console-standalone-6.1.3.jar execute -cp out --scan-class-path
\end{lstlisting}
% +++++++++++++++++++++++++++++++++++++++ A lst


	\small With a build tool, one command compiles and runs every test: \hRemark{Not run here}


% +++++++++++++++++++++++++++++++++++++++ V lst
%: -Lst. maven and gradle
\begin{lstlisting}[style=codesnippet, language={}, basicstyle=\ttfamily\scriptsize]
mvn test          pom.xml: org.junit.jupiter:junit-jupiter:6.1.3, scope test
gradle test       build.gradle: testImplementation 'org.junit.jupiter:junit-jupiter:6.1.3'
\end{lstlisting}
% +++++++++++++++++++++++++++++++++++++++ A lst


	\footnotesize JUnit 6.1.3 is the current stable version (November 2026; it needs Java 17); the labs use JUnit 5 from Eclipse. The test
	code is the same. This deck's own sessions use a small launcher, \hCode{RunTests.java}, with the jars of Eclipse.
\end{frame}
% ~~~~~~~~~~~~~~~~~~~~~~~~~~~~~~~~~~~~~~ A frm




% ~~~~~~~~~~~~~~~~~~~~~~~~~~~~~~~~~~~~~~ V
\hSection[End]{Summary}
% ~~~~~~~~~~~~~~~~~~~~~~~~~~~~~~~~~~~~~~ A




% ~~~~~~~~~~~~~~~~~~~~~~~~~~~~~~~~~~~~~~ V frm
\begin{frame}[t,fragile] \frametitle{Cheat sheet}


% +++++++++++++++++++++++++++++++++++++++ V lst
%: -Lst. cheat sheet
\begin{lstlisting}[style=codesnippet, language=Java, basicstyle=\ttfamily\scriptsize]
class MyException extends Exception { MyException(String m, Throwable c) { super(m, c); } }
this.owner = Objects.requireNonNull(owner, "owner");          // null check in one line
catch (NumberFormatException e) { throw new BadRecordException("record 3", e); }  // keep the cause
catch (NumberFormatException | ArrayIndexOutOfBoundsException e) { ... }        // multi-catch
try (Scanner in = new Scanner(f); PrintWriter out = new PrintWriter(g)) { ... }   // two resources
log.log(Level.WARNING, "no scores file", e);                  // log, with the trace

@BeforeEach void setUp() { ... }                              // a fresh fixture for every test
X e = assertThrows(X.class, () -> call());  assertEquals("msg", e.getMessage());
assertAll(() -> assertEquals(a, b), () -> assertEquals(c, d));
@ParameterizedTest @CsvSource({"60, D", "59, F"}) void t(int s, char g) { ... }
@Timeout(1) @Disabled("why") @DisplayName("a sentence")
\end{lstlisting}
% +++++++++++++++++++++++++++++++++++++++ A lst


\end{frame}
% ~~~~~~~~~~~~~~~~~~~~~~~~~~~~~~~~~~~~~~ A frm




% ~~~~~~~~~~~~~~~~~~~~~~~~~~~~~~~~~~~~~~ V frm
\begin{frame}[t,fragile] \frametitle{Projects}
	\begin{enumerate}\small
		\item \hCode{Account}: add \hCode{transfer(Account to, int amount)}. What if it fails halfway? Write the tests first, including the
			  balances after a failed transfer
		\item \hCode{Records}: read a real file; collect all bad records with their causes instead of stopping at the first
		\item A \hCode{Grades} library: parameterized tests for every boundary, a \hCode{@BeforeEach} fixture of a class list,
			  a checked \hCode{GradeFileException} for unreadable files
		\item Turn a lab of this term into a Maven or Gradle project and run its tests with one command
		\item Code review: find three \hCode{catch} blocks in your old labs that swallow or print-and-go-on. Fix them
	\end{enumerate}
\end{frame}
% ~~~~~~~~~~~~~~~~~~~~~~~~~~~~~~~~~~~~~~ A frm




% ~~~~~~~~~~~~~~~~~~~~~~~~~~~~~~~~~~~~~~ V frm
\begin{frame}[t,fragile] \frametitle{References}
	\small The Java Tutorials, lesson Exceptions \cite{TheJavaTutorials}; JUnit User Guide \cite{homeJunit}; Big Java \cite{horstmann2016big}.

	\vspace{1em}
	\setbeamertemplate{bibliography item}[text]
	\bibliographystyle{IEEEtran}
	\tiny\bibliography{../../hblib/hblibDataStructures}
\end{frame}
% ~~~~~~~~~~~~~~~~~~~~~~~~~~~~~~~~~~~~~~ A frm




% ~~~~~~~~~~~~~~~~~~~~~~~~~~~~~~~~~~~~~~ V frm
\begin{frame}[t,fragile]
	\begin{center}
		\vspace{4cm}
		\hRemark{\Huge Questions?}
	\end{center}
\end{frame}
% ~~~~~~~~~~~~~~~~~~~~~~~~~~~~~~~~~~~~~~ A frm



\end{document}
